\chapter{DCQCN: Congestion Control for Large-Scale RDMA Deployments}
\paperinfo{Y. Zhu, H. Eran, D. Firestone, C. Guo, M. Lipshteyn, Y. Liron, J. Padhye, S. Raindel, M. Haj Yahia, M. Zhang (Microsoft, Mellanox, UC Santa Barbara)}{ACM SIGCOMM 2015}

\begin{keybox}[The paper in three sentences]
To run RDMA over its ordinary IP/Ethernet datacenter (RoCEv2), Microsoft needs a network that never drops packets, and the tool for that,
PFC, does it with a blunt per-link ``PAUSE'' that makes flows unfair and spreads congestion to flows that are not even involved. DCQCN
adds per-flow congestion control on the NICs: switches ECN-mark packets, the receiver NIC turns marks into Congestion Notification
Packets (CNPs), and the sender NIC cuts and then re-grows its \emph{rate}, so queues rarely get long enough to trigger PFC at all. The paper also
works out how to set switch buffer thresholds so that ECN always fires before PFC, and uses a fluid model to choose parameters that make
flows converge to a fair share.
\end{keybox}

% ------------------------------------------------------------------
\section{Why RDMA, and why it needs a lossless network}

Microsoft wanted two things that TCP could not give. The first was 40\,Gb/s for storage traffic without burning CPU: at 40\,Gb/s, TCP
used over 20\% of all cores just moving data, and in a cloud a core spent on TCP is a core that cannot be sold as a VM. The second was
latency of a few microseconds for caches and machine learning: a 2\,KB TCP transfer took 25.4\,$\mu$s, versus 1.7\,$\mu$s with RDMA.

With \key{RDMA}, the application registers a memory buffer with the NIC, and the NIC itself moves data directly between that memory and the
remote machine's memory. The whole transport protocol runs in NIC hardware, so the CPU is almost idle (under 3\% for the RDMA client) and the
kernel is bypassed. To keep that hardware simple, RDMA's transport assumes \textbf{the network does not drop packets}. Its loss recovery is
crude (go-back-N, see Section~\ref{sec:dqpractical}), and it falls apart if losses become common.

RDMA traditionally ran on InfiniBand, a separate network technology. Microsoft runs IP and Ethernet, and will not build a second network.
\key{RoCEv2} keeps InfiniBand's transport layer but carries it in \textbf{UDP/IP over Ethernet}: IP so that packets can be routed, and UDP so
that switches can spread flows over equal-cost paths (ECMP) by hashing the UDP ports. Ethernet, however, drops packets when buffers fill. To
make it lossless, RoCEv2 relies on \key{PFC}.

% ------------------------------------------------------------------
\section{How PFC works, and how it goes wrong}

\subsection{The mechanism}
Every switch (and NIC) watches its \emph{ingress} queues, which hold packets that came in on a given port. When one grows past a threshold, the
switch sends a \textbf{PAUSE} frame back on that port to whoever is directly upstream (another switch or a host NIC). The upstream device
stops transmitting on that link until it receives a RESUME, which is sent when the queue has drained. A switch never drops for lack of buffer,
because it stops the sender before the buffer runs out. PFC has 8 priority classes, and a PAUSE applies to one class on one link.

The key weakness: \textbf{PAUSE stops a whole link (for one class), not a flow.} Every flow on that link stops, whether or not it is causing
the congestion. Two experiments on a 3-tier Clos testbed (ToRs T1--T4, leaves L1--L4, spines S1--S2, all 40\,Gb/s) show what that does.

\subsection{Problem 1: unfairness}
\paperfig[0.42]{dq_fig3_unfair}{PFC unfairness. H1--H4 all send to R, so the link T4$\to$R is the bottleneck. H4's flow enters T4 alone on
port P2, while H1--H3 share ports P3 and P4 (as ECMP happens to hash them).}

Four hosts send to R, so each should get 10\,Gb/s. T4's buffer fills, and T4 pauses its incoming ports P2, P3 and P4 in turn, so each
\emph{port} gets roughly an equal share of the time. But P2 carries only H4's flow, while P3 and P4 carry H1, H2 and H3 between them. H4 gets
a whole port's share to itself, and the other three split the remaining two. Measured: H4 got up to 20\,Gb/s, and H4's \emph{worst} run was better
than the \emph{best} run of any of H1--H3. PFC divides capacity between ports, not between flows.

\subsection{Problem 2: victim flows (congestion spreading)}
\paperfig[0.42]{dq_fig4_victim}{The victim-flow problem. H11--H14 (under T1) send to R (under T4). The flow VS$\to$VR shares only T1's uplinks
with them and is not going to the congested link, yet its throughput collapses.}

H11--H14 all send to R. T4$\to$R is congested, so T4 pauses its incoming links from L3 and L4. Now \emph{their} queues for T4 fill, so they pause
\emph{their} incoming links from S1 and S2. The spines fill and pause L1 and L2, and L1 and L2 pause T1's uplinks. T1 then has to stop its own
senders. The pause has climbed back, hop by hop, across the whole network.

The victim flow VS$\to$VR also leaves T1 over those uplinks. It is going somewhere uncongested, but it sits behind the paused packets and
stops too (head-of-line blocking). It should get 20\,Gb/s; it got 10. When two more senders (H31, H32, under T3) also sent to R, through paths
sharing \emph{no link at all} with the victim, they lengthened the pauses at L3/L4 and the spines, and the victim fell to 4.5\,Gb/s.

\subsection{Why the obvious fixes do not work}
\begin{itemize}
  \item \textbf{ECMP} was already on in these experiments, and it did not help.
  \item \textbf{More PFC priorities}: there are only 8, the experiments can be made worse just by adding senders, and flows inside one class still
        hit the same problem.
  \item \textbf{QCN} (an IEEE standard for exactly this) sends congestion feedback from the switch directly to the source, which it identifies by the
        packet's \emph{Ethernet} addresses. In an IP-routed network the Ethernet header is rewritten at every router hop, so the switch does not know
        who the source is. Fixing that needs new switch ASICs, which takes months to years.
  \item \textbf{DCTCP / iWarp / TCP-Bolt} have slow start, which hurts bursty storage traffic, or run in host software (CPU cost), or put a full TCP
        stack in the NIC (complex, and slow to recover from loss).
\end{itemize}
The real fix is \textbf{per-flow congestion control}: slow each flow down \emph{before} the queues get long enough to trigger PFC. PFC then
stays as a rarely used safety net.

% ------------------------------------------------------------------
\section{What DCQCN had to satisfy}
\begin{enumerate}
  \item Work on \textbf{L3-routed, lossless} networks.
  \item \textbf{Run on the NIC}, with low CPU cost.
  \item \textbf{Hyper-fast start}: most transfers are small and the network is usually idle, so a new flow should start at full speed, not ramp up
        slowly.
  \item Converge quickly to a fair share, keep queues short and stable, and keep links fully used.
  \item \textbf{No custom switch features}: only what every switch already has (RED and ECN marking).
\end{enumerate}

% ------------------------------------------------------------------
\section{The DCQCN control loop}

DCQCN borrows its marking and its $\alpha$ estimate from DCTCP, and its rate-increase machinery from QCN. It has three roles (names from QCN):
\begin{center}
\begin{tabular}{@{}lll@{}}
\toprule
Role & Runs on & Job\\
\midrule
\key{Congestion Point (CP)} & switch egress queue & ECN-mark packets when the queue is long\\
\key{Notification Point (NP)} & receiver NIC & turn ECN marks into CNPs back to the sender\\
\key{Reaction Point (RP)} & sender NIC & cut the sending rate on a CNP, re-grow it otherwise\\
\bottomrule
\end{tabular}
\end{center}

\textbf{One congestion event, end to end:}
\begin{enumerate}
  \item A flow is sending at line rate. Others join, and the queue at a switch port starts to grow.
  \item \textbf{CP:} the switch starts ECN-marking packets leaving that queue (more marks the longer the queue).
  \item \textbf{NP:} the receiver's NIC sees a marked packet and sends a CNP to the sender, at most one per 50\,$\mu$s for that flow.
  \item \textbf{RP:} the sender's NIC receives the CNP and cuts the flow's rate. How hard it cuts depends on how much congestion it has
        seen recently ($\alpha$).
  \item The queue drains. No more marks, no more CNPs.
  \item \textbf{RP:} with no CNPs arriving, the sender raises its rate in steps driven by a timer and a byte counter: first quickly back towards
        the old rate, then slowly probing above it.
  \item If the queue builds again, marks resume and the cycle repeats.
\end{enumerate}
The rest of this section explains each step.

\subsection{Step 1: the switch marks (CP)}
\paperfig[0.36]{dq_fig5_red}{The switch's marking curve: no marking below $K_{min}$, probability rising linearly to $P_{max}$ at $K_{max}$,
everything marked above $K_{max}$.}
When a packet leaves an egress queue of length $q$, the switch ECN-marks it with probability
\[
  p(q)=\begin{cases}0 & q\le K_{min}\\[2pt] \dfrac{q-K_{min}}{K_{max}-K_{min}}\,P_{max} & K_{min}<q\le K_{max}\\[6pt] 1 & q>K_{max}\end{cases}
\]
This is plain RED, which every switch supports. DCTCP uses a ``cliff'' ($K_{min}=K_{max}=K$, $P_{max}=1$: mark everything above $K$).
DCQCN ends up using a long, gentle slope instead ($K_{min}$ = 5\,KB, $K_{max}$ = 200\,KB, $P_{max}$ = 1\%). Section~\ref{sec:dqparams}
explains why that matters.

\subsection{Step 2: the receiver sends CNPs (NP)}
\paperfig[0.33]{dq_fig6_np}{The NP state machine, per flow: send a CNP immediately for the first marked packet, then at most one CNP every
50\,$\mu$s, as long as some packet in that window was marked.}
DCTCP echoes every mark back in an ACK. DCQCN cannot afford that: handling marked packets and building CNPs is expensive work for the NIC.
So the receiver NIC rate-limits its feedback per flow:
\begin{itemize}
  \item a marked packet arrives and no CNP has gone to this flow in the last 50\,$\mu$s: send a CNP \textbf{now};
  \item after that, send \textbf{at most one CNP per 50\,$\mu$s}, and only if at least one packet in that window was marked.
\end{itemize}
A CNP therefore carries one bit of information: ``you were marked at least once in the last 50\,$\mu$s''. It says nothing about how
\emph{many} packets were marked. If there is no congestion, there are no marks and no CNPs at all. CNPs are sent at high priority, so they are
not stuck behind the congestion they report.

\subsection{Step 3: the sender cuts its rate (RP)}
\paperfig[0.62]{dq_fig7_rp}{The RP state machine. A CNP cuts the rate and resets all counters. Expiries of the rate timer and the byte counter
drive the increase phases. A separate alpha timer decays $\alpha$.}

For each flow it is rate-limiting, the sender NIC keeps a \textbf{current rate} $R_C$ (what it sends at), a \textbf{target rate} $R_T$ (the rate it
is trying to get back to), and \textbf{$\alpha$}, an estimate of how congested the flow has recently been, starting at 1. On a CNP:
\[
  R_T \leftarrow R_C,\qquad R_C \leftarrow R_C\Big(1-\frac{\alpha}{2}\Big),\qquad \alpha \leftarrow (1-g)\,\alpha + g
\]
In words:
\begin{itemize}
  \item \textbf{Remember where we were} ($R_T$), so we know what to climb back towards.
  \item \textbf{Cut by $\alpha/2$.} If $\alpha$ is near 1 (CNPs keep coming), the rate is halved. If $\alpha$ is small (this is a rare CNP), the cut
        is gentle. This is exactly DCTCP's idea: cut in proportion to how persistent the congestion is.
  \item \textbf{Raise $\alpha$} a little, because we just saw congestion. $g$ is the averaging weight; the deployed value is $1/256$.
\end{itemize}
It also resets the increase counters (below). Separately, every 55\,$\mu$s in which \emph{no} CNP arrives, $\alpha$ decays:
$\alpha \leftarrow (1-g)\alpha$. That interval must be longer than the 50\,$\mu$s CNP interval. Otherwise $\alpha$ would decay in the gap between two
CNPs of a flow that is in fact still congested.

\subsection{Step 4: the sender re-grows its rate (RP)}
Because silence means ``no congestion'', the sender has to decide by itself when to speed up. It has two triggers, both from QCN:
\begin{itemize}
  \item a \textbf{timer} that fires every $T$ (deployed: 55\,$\mu$s), and
  \item a \textbf{byte counter} that fires every $B$ bytes sent (deployed: 10\,MB).
\end{itemize}
The timer matters when the rate is very low: a flow cut to a trickle would take ages to send $B$ bytes, so the byte counter alone would
leave it stuck. Each firing performs one increase step. The NIC counts how many times each trigger has fired since the last CNP (counts $T$ and $BC$),
and the counts decide which kind of step is taken, with $F = 5$:
\begin{center}
\begin{tabularx}{\linewidth}{@{}lllX@{}}
\toprule
Phase & When & Update & Purpose\\
\midrule
\key{Fast recovery} & both counts $< 5$ & $R_C \leftarrow (R_T+R_C)/2$ & close half the gap back to the rate before the cut, each step\\
\key{Additive increase} & one count $\ge 5$ & $R_T \leftarrow R_T + R_{AI}$, then $R_C \leftarrow (R_T+R_C)/2$ & probe gently above the old rate ($R_{AI}=40$\,Mb/s)\\
\key{Hyper increase} & both counts $> 5$ & larger target steps (as in QCN) & a long quiet period: ramp up fast\\
\bottomrule
\end{tabularx}
\end{center}

\textbf{Worked example.} A flow is at $R_C = 20$\,Gb/s with $\alpha = 0.5$ ($g = 1/256$). One CNP arrives, then no more.
\begin{center}
\small
\begin{tabular}{@{}llccc@{}}
\toprule
Event & What happens & $R_T$ & $R_C$ (Gb/s) & $\alpha$\\
\midrule
CNP & $R_T{=}20$; $R_C = 20(1-0.25)$; $\alpha$ nudged up & 20 & 15 & 0.502\\
increase 1 & fast recovery: $(20+15)/2$ & 20 & 17.5 & \\
increase 2 & fast recovery & 20 & 18.75 & \\
increase 3 & fast recovery & 20 & 19.38 & \\
increase 4 & fast recovery & 20 & 19.69 & \\
increase 5 & fast recovery & 20 & 19.84 & \\
increase 6 & additive: $R_T$ up 0.04, then average & 20.04 & 19.94 & \\
every 55\,$\mu$s & no CNP: $\alpha$ decays & & & slowly falls\\
\bottomrule
\end{tabular}
\end{center}
The shape is: a cut sized by recent congestion, a quick return to just below the old rate, then a slow climb above it. If the climb
causes a queue again, a CNP arrives and the cycle repeats around the fair rate.

\subsection{No slow start}
A new flow starts at \textbf{full line rate} (or at the local QoS rate if the host already has active flows). Most transfers are small and the network
is usually idle, so this makes the common case as fast as possible. It is also why \textbf{PFC is still needed}: a burst of line-rate starts can
overflow a buffer before any CNP has had time to arrive, and PFC catches that.

\subsection{Why a rate and not a window, and the hardware cost}
A window-based design (like TCP) needs per-ACK bookkeeping. A rate is easier to enforce in hardware: the NIC already has rate limiters that pace
every packet, so DCQCN only adjusts their settings. Because the rate does not depend on measuring RTTs, DCQCN is not confused by congestion
on the reverse path. Per rate-limited flow, the NIC keeps one timer, one counter, $\alpha$ and a few rates, all on the NIC chip. The ConnectX-3
Pro can build one CNP every 1--5\,$\mu$s. At 40\,Gb/s about 166 full-size packets arrive per 50\,$\mu$s, so it can serve CNPs for 10--20
congested flows at once (ConnectX-4: over 200).

% ------------------------------------------------------------------
\section{Switch buffer settings: making ECN act before PFC}

DCQCN only helps if switches mark packets (and senders slow down) \emph{before} PFC pauses anything. If PFC fires first, all the
unfairness and spreading come back. There are three thresholds on the switch (Arista 7050QX32: 32 ports of 40\,Gb/s, a 12\,MB buffer shared by all
ports, 8 priorities):
\begin{itemize}
  \item \textbf{Headroom $t_{flight}$.} After a switch sends PAUSE, packets keep arriving for a while: those already on the wire, plus whatever the
        upstream device sends before it processes the PAUSE (it cannot stop in the middle of a packet). Buffer must be reserved for them on every
        port and priority, or they would be dropped. With a 1500-byte MTU this comes to \textbf{22.4\,KB} per port per priority.
  \item \textbf{PFC threshold $t_{PFC}$}, on an \emph{ingress} queue: PAUSE is sent when the queue exceeds it. Whatever is left after the headroom,
        divided over all ports and priorities, gives at most \textbf{24.47\,KB}. RESUME is sent when the queue falls two MTUs below it.
  \item \textbf{ECN threshold $t_{ECN}$} ($= K_{min}$), on an \emph{egress} queue: marking starts above it.
\end{itemize}

\textbf{The catch: PFC watches ingress queues, ECN watches egress queues.} Worst case: everything waiting in the switch's 32 egress queues
arrived through \emph{one} ingress port. That ingress queue then holds the sum of all of them. For ECN to start on at least one egress queue
before this one ingress queue reaches $t_{PFC}$, we need $32 \times t_{ECN} < t_{PFC}$, which forces $t_{ECN}$ below one kilobyte (the paper
gets 0.85\,KB). That is smaller than a single packet, so it is impossible.

\textbf{The fix: a dynamic PFC threshold.} The buffer is shared, so a fixed per-queue limit wastes it. When most of the buffer is empty,
there is no hurry to pause anyone. The switch chip (Broadcom Trident II) can set the PFC threshold from the buffer that is currently free:
\[
  t_{PFC} = \frac{\beta\,(B - 8n\,t_{flight} - s)}{8}
\]
where $s$ is the buffer currently occupied. When the buffer is nearly empty, the threshold is large, and it shrinks as the buffer fills. Just before
ECN fires anywhere, the occupancy is at most $s \le n\,t_{ECN}$. Substituting that, and requiring ECN to come first, gives the paper's condition
\[
  t_{ECN} < \frac{\beta\,(B - 8n\,t_{flight})}{8n(\beta+1)},
\]
which with $\beta = 8$ gives $t_{ECN} < $ \textbf{21.75\,KB}, comfortably more than one packet. A larger $\beta$ leaves more room for $t_{ECN}$
and makes PFC fire less often.

This is a worst-case analysis (all 8 priorities in use). It does not promise that PFC never fires, only that at each switch \textbf{PFC never fires
before ECN}. Between the first mark and the sender's reaction, PFC may still pause briefly, and that is the job it is kept for.

% ------------------------------------------------------------------
\section{Choosing the parameters: why the obvious settings fail}\label{sec:dqparams}

The authors wrote a \key{fluid model}: differential equations for the queue, $\alpha$, $R_T$ and $R_C$ of $N$ flows sharing one bottleneck of capacity
$C$, with a 50\,$\mu$s feedback delay. Its equilibrium is $R_C = C/N$ (everyone gets a fair share), with a marking probability below 1\%, so the
queue settles just above $K_{min}$. The model matched the real NIC's behaviour closely, so they used it to test parameter choices. The test:
\textbf{one flow is at 40\,Gb/s and a second starts at 0. Do they converge to 20 and 20?}

\subsection{First try: QCN's and DCTCP's own defaults, and why they fail}
Byte counter $B = 150$\,KB, timer $T = 1.5$\,ms, DCTCP cliff marking at 40\,KB, $g = 1/16$. \textbf{The flows never converge.}

The reason is the byte counter. With $T = 1.5$\,ms the timer hardly ever fires, so increases are driven by bytes sent, and a faster flow sends
bytes faster:
\begin{center}
\begin{tabular}{@{}lll@{}}
\toprule
Flow & Time to send 150\,KB & Increase steps per ms\\
\midrule
fast flow, 30\,Gb/s & 40\,$\mu$s & 25\\
slow flow, 10\,Gb/s & 120\,$\mu$s & about 8\\
\bottomrule
\end{tabular}
\end{center}
The fast flow climbs three times as often, so it grabs back whatever it gives up and the gap never closes. QCN gets away with this because its
feedback is probabilistic per packet; DCQCN's timer-driven CNPs are not.

\subsection{Fix 1: let the timer drive the increase}
Make the timer fast, $T = 55$\,$\mu$s, so that both flows get an increase step every 55\,$\mu$s regardless of their rate. Then make the byte
counter large (10\,MB, which takes 2.7\,ms even at 30\,Gb/s), so it rarely fires and does not push flows into hyper increase too early.
Now both flows get the \emph{same} increases and proportional cuts. The faster flow loses more in absolute terms on each cut and gains the same,
so the gap shrinks. The timer cannot be shorter than the 50\,$\mu$s CNP interval, or a flow would increase between two CNPs it is still
receiving. Hence 55\,$\mu$s: 50 plus a 5\,$\mu$s margin.

\subsection{Fix 2: let bigger flows hear about congestion more often}
With cliff marking, once the queue is above $K$ \emph{every} packet is marked, so every flow gets a CNP every 50\,$\mu$s, big or small. The
feedback does not discriminate. With a gentle RED slope and a small $P_{max}$, a packet is marked only occasionally (say $p = 0.5\%$). A flow
then gets a CNP in a given 50\,$\mu$s window only if \emph{at least one} of its packets in that window was marked, and a bigger flow has more
packets in the window:
\begin{center}
\begin{tabular}{@{}llll@{}}
\toprule
Flow & 1500\,B packets per 50\,$\mu$s & P(some packet marked) $= 1-(1-0.005)^{n}$ & CNPs\\
\midrule
30\,Gb/s & 125 & $1 - 0.995^{125} \approx 0.47$ & in about half the windows\\
10\,Gb/s & 42 & $1 - 0.995^{42} \approx 0.19$ & in about 1 window in 5\\
\bottomrule
\end{tabular}
\end{center}
(The 0.5\% is illustrative; the paper only states that $p < 1\%$ at equilibrium.) The big flow is cut about 2.5 times as often, so it
backs off faster, which pushes the flows together. The deployed setting uses \textbf{both} fixes. Probabilistic marking alone made throughput
jumpy in the testbed, but the combination worked best, especially with multiple bottlenecks.

\subsection{The remaining knobs}
\begin{itemize}
  \item \textbf{$g$}: smaller $g$ makes $\alpha$ change more slowly, which gave lower and steadier queues in 2:1 to 16:1 incast, at the price of slightly
        slower convergence. Chosen: $1/256$.
  \item \textbf{$K_{min} = 5$\,KB} looks very shallow, but the marking probability just above it is tiny, so the queue actually settles about ten times
        higher, which is enough to keep the link 100\% busy. DCTCP needs a deep threshold (around 160\,KB at 40\,Gb/s) because software sending is bursty;
        DCQCN's hardware pacing is smooth, so a shallow queue works. In a 20:1 incast, the 95th-percentile queue was 76.6\,KB for DCQCN against
        162.9\,KB for DCTCP.
  \item \textbf{$R_{AI}$ and $F$} trade convergence speed against oscillation. With the chosen values there is no buffer starvation at 16:1 incast;
        halving $R_{AI}$ would handle 32:1 but converge more slowly.
\end{itemize}

\begin{center}
\begin{tabularx}{\linewidth}{@{}llX@{}}
\toprule
Parameter & Deployed & Reason\\
\midrule
rate-increase timer & 55\,$\mu$s & increase independent of rate; just above the CNP interval\\
byte counter & 10\,MB & keep it from dominating the increase or triggering hyper increase early\\
$K_{min}$ / $K_{max}$ / $P_{max}$ & 5\,KB / 200\,KB / 1\% & long gentle slope, so bigger flows get more CNPs\\
$g$ & 1/256 & low, steady queues\\
$\alpha$ decay interval & 55\,$\mu$s & longer than the CNP interval\\
$F$, $R_{AI}$ & 5, 40\,Mb/s & convergence speed vs oscillation\\
\bottomrule
\end{tabularx}
\end{center}

With these settings, 1 to 19 senders into one receiver always kept total throughput above 39\,Gb/s with the switch queue under 100\,KB. Back in the
two experiments: the four senders split the bottleneck equally, and the victim flow kept its throughput as senders were added. In a storage
benchmark with incast, the spine switches received about 3000 PAUSE messages instead of over 6 million.

% ------------------------------------------------------------------
\section{Practical issues}\label{sec:dqpractical}
\begin{itemize}
  \item \textbf{Multiple bottlenecks.} A flow crossing two congested links gets CNPs caused by both, so it ends up below its fair share. Probabilistic
        marking reduces this effect but does not remove it.
  \item \textbf{Packet loss still happens}, from corruption, bad cables and misconfiguration; PFC only prevents overflow drops. RoCEv2 recovers with
        go-back-N: the receiver spots a gap in sequence numbers, and everything from the lost packet onward is resent. That works only while loss stays
        rare, and throughput falls sharply above about 0.1\% loss.
  \item \textbf{Deadlock.} PAUSEs can deadlock if buffers wait on each other in a cycle. In a Clos datacenter, traffic never goes through a server or
        more than two ToRs, so no cycle can form unless something is broken or misconfigured.
  \item \textbf{Coexisting with TCP} is handled by giving RDMA and TCP different priority classes, each with its own buffers and limits. TCP-friendliness
        was not a goal.
\end{itemize}

\begin{keybox}[The big idea]
Congestion control always trades responsiveness against stability. DCQCN pairs a \textbf{fast but blunt} mechanism (PFC, which stops loss
instantly) with a \textbf{precise but slower} one (per-flow rate control on the NIC, which keeps queues short so PFC rarely fires). Together they are
lossless in the short term and fair and stable in the long term.
\end{keybox}

\section{DCTCP vs.\ QCN vs.\ DCQCN}
\begin{center}
\small
\begin{tabularx}{\linewidth}{@{}lXXX@{}}
\toprule
 & \textbf{DCTCP} & \textbf{QCN} & \textbf{DCQCN}\\
\midrule
Network & IP, lossy & Ethernet L2 only & IP (RoCEv2), lossless via PFC\\
Runs in & host kernel & switch + NIC & NIC; switch only does standard RED/ECN\\
Controls & window & rate & rate (hardware pacing)\\
Feedback & ECN echoed in every ACK & quantized value, switch $\to$ source & CNP from receiver, $\le$ 1 per 50\,$\mu$s\\
Cut & $W(1-\alpha/2)$ & by feedback value & $R_C(1-\alpha/2)$\\
Growth & slow start + additive & fast recovery / additive / hyper & same as QCN, no slow start\\
\bottomrule
\end{tabularx}
\end{center}

\begin{keybox}[Exam one-liner]
DCQCN makes RoCEv2 RDMA work at datacenter scale by adding per-flow, NIC-based rate control on top of PFC. Switches RED/ECN-mark,
the receiver NIC sends at most one CNP per 50\,$\mu$s per flow, and the sender NIC cuts $R_C$ by $\alpha/2$ (remembering $R_T$), then recovers
with timer-driven fast recovery, additive and hyper increase, with no slow start. A dynamic PFC threshold ($\beta = 8$) guarantees that ECN fires
before PFC. A fluid model shows that a fast timer (55\,$\mu$s) and gentle probabilistic marking ($P_{max}$ = 1\%) are needed for flows to converge.
\end{keybox}

\section{Revision checklist}
\begin{itemize}
  \item Why RDMA (CPU and latency); RoCEv2 = IB transport over UDP/IP/Ethernet; needs lossless, hence PFC.
  \item PFC pauses a link, not a flow: unfairness (H4 got a whole port) and victim flows (pauses climb hop by hop).
  \item QCN fails on IP networks (Ethernet header rewritten per hop); real fix = per-flow congestion control.
  \item Loop: CP marks (RED) $\to$ NP sends CNP ($\le$ 1 per 50\,$\mu$s) $\to$ RP cuts $R_C(1-\alpha/2)$, $R_T = $ old rate, $\alpha$ up; $\alpha$ decays every 55\,$\mu$s without CNP.
  \item Increase: timer + byte counter; fast recovery halves the gap, additive adds 40\,Mb/s to target, hyper after long quiet; no slow start.
  \item Buffers: headroom 22.4\,KB; static PFC limit forces $t_{ECN} < 1$ packet (impossible); dynamic $t_{PFC}$ with $\beta = 8$ gives $t_{ECN} < 21.75$\,KB.
  \item Convergence: byte-counter-driven increase favours fast flows (25 vs 8 steps/ms); fix with 55\,$\mu$s timer and RED slope with small $P_{max}$ (big flows get more CNPs).
  \item Practical: parking lot mitigated, go-back-N needs loss $<$ 0.1\%, no deadlock in Clos, TCP in a separate class.
\end{itemize}
